\pdfbookmark[0]{Abstract}{Abstract}

\begin{otherlanguage}{american}
	\chapter*{Abstract}
Software quality is often difficult for decision-makers to grasp, especially when codebases and metrics are unfamiliar. Visual representations can enhance the interpretability of technical analyses. However, there is a gap between research prototypes and practice-ready visualizations. This thesis addresses that gap through a systematic comparison, a unified implementation, and targeted improvements to practice-relevant visualizations.

To compile a comprehensive inventory of current software visualizations, 26 approaches from research and practice are reviewed and assessed. Based on this, the five best visualization types are selected for a unified implementation. Treemap, Streetmap, CodeCity, Sunburst, and Circular Treemap visualizations are examined in a consistent configuration, each encoding three metrics via area, height, and color.

A core focus is the technical improvement of treemap visualization based on the Squarify algorithm. In existing implementations, nodes frequently disappear—an issue rarely addressed in research or practice. A two-stage algorithm is proposed and its quality evaluated on a broad open-source dataset using criteria such as node visibility and area proportionality.

A realistic user study with scenario-based assessments and a direct expert evaluation then validate the findings of the inventory and the treemap optimization. Across the literature review, unified implementations, and the user survey, treemaps most strongly support the interpretability of quality statements. Circular treemaps also perform well on structure-oriented tasks and are popular with users. The improved treemap reduces the number of missing nodes by up to a factor of 4, improves node aspect ratios to a median of approximately 2{.}5, and halves the deviation in area proportionality compared to the baseline. The work thus provides a practice-ready algorithm and robust guidelines for application in audits, health checks, and general communication about software.

The grounded comparison of the 26 approaches further shows which visualizations are particularly suitable for presenting software quality and structure to non-technical decision-makers, where the main opportunities for improvement lie, and how they can be addressed. Overall, the thesis establishes an evidence-based framework—combining literature review, unified implementations, and an empirical user survey—that supports the selection of appropriate visualizations in practice.
\end{otherlanguage}